\documentclass[Orbiter User Manual.tex]{subfiles}

\begin{document}

\section{Selected constants and parameters}
\label{sec:constants}

These data may be useful when planning missions across the solar system.


\subsection{Astrodynamic constants}

%\begin{table}[H]
	%\centering
	\begin{longtable}{ p{0.35\textwidth} p{0.14\textwidth} p{0.42\textwidth} }
	\hline
	\textbf{Constant} & \textbf{Symbol} & \textbf{Value}\\
	\hline
	Julian day & d & 86400\,s\\
	\hline
	Julian year & yr & 365.25\,d\\
	\hline
	Julian century & Cy & 36525\,d\\
	\hline
	Speed of light & c & 299792458\,m/s\\
	\hline
	Gaussian gravitational constant & k &
		0.01720209895 (AU\textsup{3}/d\textsup{2})\textsup{1/2} \\*
	\hline
	\caption{Defining constants}
	\end{longtable}
%\end{table}

%TODO add vesta, ceres, etc
%\begin{table}[H]
	%\centering
	\begin{longtable}{ p{0.35\textwidth} p{0.14\textwidth} p{0.42\textwidth} }
	\hline
	\textbf{Constant} & \textbf{Symbol} & \textbf{Value}\\
	\hline
	Mean sidereal day & & 86164.09054\,s = 23:56:04.09054\\
	\hline
	Sidereal year (quasar ref. frame) & & 365.25636\,d\\
	\hline
	Light time for 1\,AU & $\tau_{A}$ & 499.004783806 ($\pm$\,0.00000001)\,s\\
	\hline
	Gravitational constant & G & 6.67259 ($\pm$\,0.00030) $\cdot$
		10\textsup{-11}\,kg\textsup{-1}\,m\textsup{3}\,s\textsup{-2}\\
	\hline
	General precession in longitude & & 5028.83 ($\pm\,$0.04)\,arcsec/Cy\\
	\hline
	Obliquity of ecliptic (J2000.0) & $\epsilon$ &
		84381.412 ($\pm\,$0.005)\,arcsec\\
	\hline
	Mass: Sun / Mercury & & 6023600. ($\pm$\,250.)\\
	\hline
	Mass: Sun / Venus & & 408523.71 ($\pm\,$0.06)\\
	\hline
	Mass: Sun / (Earth + Moon) & & 328900.56 ($\pm$\,0.02)\\
	\hline
	Mass: Sun / (Mars system) & & 3098708. ($\pm$\,9.)\\
	\hline
	Mass: Sun / (Jupiter system) & & 1047.3486 ($\pm$\,0.0008)\\
	\hline
	Mass: Sun / (Saturn system) & & 3497.898 ($\pm$\,0.018)\\
	\hline
	Mass: Sun / (Uranus system) & & 22902.98 ($\pm\,$0.03)\\
	\hline
	Mass: Sun / (Neptune system) & & 19412.24 ($\pm$\,0.04)\\
	\hline
	Mass: Sun / (Pluto system) & & 1.35 ($\pm$\,0.07) $\cdot$ 10\textsup{8}\\
	\hline
	Mass: Moon / Earth & & 0.012300034 ($\pm$\,3 $\cdot$ 10\textsup{-9})\\
	\hline
	\caption{Primary constants}
	\end{longtable}
%\end{table}

%\begin{table}[H]
	%\centering
	\begin{longtable}{ p{0.35\textwidth} >{\RaggedRight}p{0.14\textwidth}
		p{0.42\textwidth} }
	\hline
	\textbf{Constant} & \textbf{Symbol} & \textbf{Value}\\
	\hline
	Astronomical unit distance & c $\tau_{A}$ = AU &
		1.49597870691 $\cdot$ 10\textsup{11} ($\pm$\,3)\,m\\
	\hline
	Heliocentric gravitational constant &
		k\textsup{2}\,AU\textsup{3}\,d\textsup{-2} = GM$_\mathrm{Sun}$ &
		1.32712440018 $\cdot$ 10\textsup{20} ($\pm$\,8 $\cdot$
			10\textsup{9})\,m\textsup{3}\,s\textsup{-2}\\
	\hline
	Mass: Earth / Moon & & 81.30059 ($\pm$\,0.00001)\\
	\hline
	\caption{Derived constants}
	\end{longtable}
%\end{table}

All values from \cite{standish1995}

\underline{Notes:}\\
Data are from the 1994 IAU file of current best estimates. Planetary ranging
determines the Earth/Moon mass ratio. The value for 1 AU is taken from JPL's
current planetary ephemeris DE-405.

\paragraph{Planetary mean orbital elements and centennial rates}

%TODO add vesta, ceres, etc
%\begin{table}[H]
	%\centering
	\begin{longtable}{ p{0.1\textwidth} p{0.13\textwidth} p{0.13\textwidth}
		p{0.1\textwidth} p{0.12\textwidth} p{0.11\textwidth} p{0.11\textwidth} }
	\hline
	\textbf{Planet} & \textbf{$a$ [AU]} & \textbf{$e$} & \textbf{$i$ [deg]} &
		\textbf{$\Omega$ [deg]} & \textbf{$\bar{\omega}$ [deg]} &
		\textbf{$L$ [deg]}\\
	\hline
	Mercury & 0.38709893 & 0.20563069 & 7.00487 & 48.33167 & 77.45645 &
		252.25084\\
	\hline
	Venus & 0.72333199 & 0.00677323 & 3.39471 & 76.68069 & 131.53298 &
		181.97973\\
	\hline
	Earth & 1.00000011 & 0.01671022 & 0.00005 & -11.26064 & 102.94719 &
		100.46435\\
	\hline
	Mars & 1.52366231 & 0.09341233 & 1.85061 & 49.57854 & 336.04084 &
		355.45332\\
	\hline
	Jupiter & 5.20336301 & 0.04839266 & 1.30530 & 100.55615 & 14.75385 &
		34.40438\\
	\hline
	Saturn & 9.53707032 & 0.05415060 & 2.48446 & 113.71504 & 92.43194 &
		49.94432\\
	\hline
	Uranus & 19.19126393 & 0.04716771 & 0.76986 & 74.22988 & 170.96424 &
		313.23218\\
	\hline
	Neptune & 30.06896348 & 0.00858587 & 1.76917 & 131.72169 & 44.97135 &
		304.88003\\
	\hline
	Pluto & 39.48168677 & 0.24880766 & 17.14175 & 110.30347 & 224.06676 &
		238.92881\\
	\hline
	\end{longtable}
%\end{table}

%TODO add vesta, ceres, etc
%\begin{table}[H]
	%\centering
	\begin{longtable}{ p{0.1\textwidth} p{0.13\textwidth} p{0.13\textwidth}
		p{0.08\textwidth} p{0.11\textwidth} p{0.1\textwidth} p{0.15\textwidth} }
	\hline
	\textbf{Planet} & \textbf{$\diffd a/\diffd t$ [AU/Cy]} &
		\textbf{$\diffd e/\diffd t$ [/Cy]} &
		\textbf{$\diffd i/\diffd t$ [$''$/Cy]} &
		\textbf{$\diffd\Omega/\diffd t$ [$''$/Cy]} &
		\textbf{$\diffd \bar{\omega}/\diffd t$ [$''$/Cy]} &
		\textbf{$\diffd L/\diffd t$ [$''$/Cy]}\\
	\hline
	Mercury & 0.00000066 & 0.00002527 & -23.51 & -446.30 & 573.57 &
		538101628.29\\
	\hline
	Venus & 0.00000092 & -0.00004938 & -2.86 & -996.89 & -108.80 &
		210664136.06\\
	\hline
	Earth & -0.00000005 & -0.00003804 & -46.94 & -18228.25 & 1198.28 &
		129597740.63\\
	\hline
	Mars & -0.00007221 & 0.00011902 & -25.47 & -1020.19 & 1560.78 &
		68905103.78\\
	\hline
	Jupiter & 0.00060737 & -0.00012880 & -4.15 & 1217.17 & 839.93 &
		10925078.35 \\
	\hline
	Saturn & -0.00301530 & -0.00036762 & 6.11 & -1591.05 & -1948.89 &
		4401052.95\\
	\hline
	Uranus & 0.00152025 & -0.00019150 & -2.09 & -1681.40 & 1312.56 &
		1542547.79\\
	\hline
	Neptune & -0.00125196 & 0.0000251 & -3.64 & -151.25 & -844.43 & 786449.21\\
	\hline
	Pluto & -0.00076912 & 0.00006465 & 11.07 & -37.33 & -132.25 & 522747.90\\
	\hline
	\end{longtable}
%\end{table}

$a$: semi-major axis, $e$: eccentricity, $i$: inclination, $\Omega$: longitude
of the ascending node, $\bar{\omega}$: longitude of perihelion, $L$: mean
longitude, Cy: Julian century, $''$: arcsec\\
Epoch = J2000 = 2000 January 1 12h\\
All values from \cite{seidelmann1992p316}

\underline{Notes:}\\
The tables contain mean orbit solutions from a 250 yr. least squares fit of the
DE 200 planetary ephemeris to a Keplerian orbit where each element is allowed to
vary linearly with time. This solution fits the terrestrial planet orbits to
$\sim 25''$ or better, but achieves only $\sim 600''$ for Saturn. Elements are
referenced to mean ecliptic and equinox of J2000 at the J2000 epoch (2451545.0
JD).


\subsection{Selected physical parameters}

%TODO add vesta, ceres, etc
%\begin{table}[H]
	%\centering
	\begin{longtable}{ p{0.1\textwidth} >{\RaggedRight}p{0.16\textwidth}
		>{\RaggedRight}p{0.15\textwidth} >{\RaggedRight}p{0.17\textwidth}
		>{\RaggedRight}p{0.13\textwidth} >{\RaggedRight}p{0.12\textwidth} }
	\hline
	\textbf{Planet} & \textbf{mean radius [km]} &
		\textbf{mass [10\textsup{23}\,kg]} &
		\textbf{density [g/cm\textsup{3}]} &
		\textbf{sidereal rotation period [h]} &
		\textbf{sidereal orbit period [yr]}\\
	\hline
	Mercury & 2440.\,$\pm$\,1. & 3.301880 & 5.427 & 1407.509 & 0.2408445\\
	\hline
	Venus & 6051.84\,$\pm$\,0.01 & 48.6855374 & 5.204 & -5832.444 & 0.6151826\\
	\hline
	Earth & 6371.01\,$\pm$\,0.02 & 59.73698968 & 5.515 & 23.93419** &
		0.9999786\\
	\hline
	Mars & 3389.92\,$\pm$\,0.04 & 6.418542 & 3.9335 $\pm$ 0.0004 & 24.622962 &
		1.88071105\\
	\hline
	Jupiter & 69911.\,$\pm$\,6. & 18986.111 & 1.326 & 9.92425 & 11.856523\\
	\hline
	Saturn & 58232.\,$\pm$\,6. & 5684.6272 & 0.6873 & 10.65622 & 29.423519\\
	\hline
	Uranus & 25362.\,$\pm$\,12. & 868.32054 & 1.318 & 17.24\,$\pm$\,0.01 &
		83.747407\\
	\hline
	Neptune & 24624.\,$\pm$\,21. & 1024.569 & 1.638 & 16.11\,$\pm$\,0.01 &
		163.72321\\
	\hline
	Pluto* & 1151 & 0.15 & 1.1 & 153.28 & 248.0208\\
	\hline
	\end{longtable}
%\end{table}

%TODO add vesta, ceres, etc
%\begin{table}[H]
	%\centering
	\begin{longtable}{ p{0.1\textwidth} >{\RaggedRight}p{0.17\textwidth}
		>{\RaggedRight}p{0.17\textwidth} >{\RaggedRight}p{0.21\textwidth}
		>{\RaggedRight}p{0.21\textwidth} }
	\hline
	\textbf{Planet} & \textbf{V(1,0) [mag.]} & \textbf{Geometric albedo} &
		\textbf{Equatorial gravity [m/s$^{2}$]} &
		\textbf{Escape velocity [km/s]}\\
	\hline
	Mercury & -0.42 & 0.106 & 3.701 & 4.435\\
	\hline
	Venus & -4.4 & 0.65 & 8.87 & 10.361\\
	\hline
	Earth & -3.86 & 0.367 & 9.780327 & 11.186\\
	\hline
	Mars & -1.52 & 0.15 & 3.69 & 5.027\\
	\hline
	Jupiter & -9.4 & 0.52 & 23.12\,$\pm$\,0.01 & 59.5\\
	\hline
	Saturn & -8.88 & 0.47 & 8.96\,$\pm$\,0.01 & 35.5\\
	\hline
	Uranus & -7.19 & 0.51 & 8.69\,$\pm$\,0.01 & 21.3\\
	\hline
	Neptune & -6.87 & 0.41 & 11.00\,$\pm$\,0.05 & 23.5\\ 
	\hline
	Pluto* & -1.0* & 0.3* & 0.655 & 1.3\\
	\hline
	\end{longtable}
%\end{table}

All values from \cite{yoder1995} except Pluto data from
\cite{seidelmann1992p706}. Mercury to Neptune masses derived from standard
gravitational parameter (GM) data in \cite{yoder1995}.\\
** Orbiter uses 23.93447 h (=23 h 56 m 4.09 s) which gives better long-term
stability.

\end{document}
